\documentclass[11pt]{article}
\usepackage[margin=1in]{geometry}
\usepackage{enumitem}
% =============================================================================
% Preambles/header.tex — shared LaTeX/Beamer preamble for this project
%
% Use from a Beamer lecture by prepending:
%     \input{header}
% and compiling with TEXINPUTS=../Preambles:$TEXINPUTS (the /compile-latex
% skill does this automatically).
%
% PALETTE CONTRACT. Color names defined here MUST match the names in
% Quarto/theme-template.scss so Beamer and Quarto renderings use the same
% palette. The scripts/check-palette-sync.sh script enforces this.
%
% To customize: edit the HEX values below AND the matching $variable in
% Quarto/theme-template.scss, then run scripts/check-palette-sync.sh.
% =============================================================================

% -----------------------------------------------------------------------------
% Engine detection + encoding
%
% Our /compile-latex skill uses XeLaTeX, where inputenc is unnecessary and
% can produce encoding warnings. Only load inputenc under pdfTeX.
% -----------------------------------------------------------------------------
\usepackage{iftex}
\ifPDFTeX
  \usepackage[utf8]{inputenc}
\fi

\usepackage{amsmath, amssymb, amsthm}
\usepackage{booktabs}
\usepackage{graphicx}

% -----------------------------------------------------------------------------
% PALETTE — mirrors Quarto/theme-template.scss variable names.
% Load xcolor and define colors BEFORE anything that references them
% (notably hyperref, hypersetup, and the Beamer color assignments).
% -----------------------------------------------------------------------------
\usepackage{xcolor}

\definecolor{primary-blue}{HTML}{012169}      % headings, accents
\definecolor{primary-gold}{HTML}{B9975B}      % emphasis, borders
\definecolor{highlight-yellow}{HTML}{F2A900}  % markers, alerts
\definecolor{light-bg}{HTML}{E8EDF5}          % title / section backgrounds
\definecolor{jet}{HTML}{1A1A1A}               % body text

% Semantic colors (used in diagrams and annotations)
\definecolor{positive}{HTML}{15803D}          % good / observed / identified
\definecolor{negative}{HTML}{B91C1C}          % bad / problematic / bias
\definecolor{neutral}{HTML}{525252}           % reference / context

% Highlight accents (match .hi-* classes in the SCSS)
\definecolor{hi-slate}{HTML}{314F4F}
\definecolor{hi-green}{HTML}{2E8B57}
\definecolor{hi-red}{HTML}{C41E3A}

% Semantic aliases so skill templates and snippets can write
% \textcolor{accent}{...} without hard-coding "primary-blue".
\colorlet{accent}{primary-blue}
\colorlet{accent2}{primary-gold}

% -----------------------------------------------------------------------------
% Hyperref — loaded AFTER xcolor + palette so link colors resolve.
% -----------------------------------------------------------------------------
\usepackage{hyperref}
\hypersetup{colorlinks=true, linkcolor=primary-blue, urlcolor=primary-blue,
            citecolor=primary-blue, pdfborder={0 0 0}}

% -----------------------------------------------------------------------------
% Course code — set once per deck (after \input{header}, before \begin{document})
% via \coursecode{CS401}. Shown in the footer on every slide so a deck's course
% is identifiable without opening the title page. Empty by default.
% -----------------------------------------------------------------------------
\makeatletter
\newcommand{\coursecode}[1]{\gdef\@coursecodetext{#1}}
\gdef\@coursecodetext{}
\makeatother

% -----------------------------------------------------------------------------
% Beamer theme assignments (only apply under Beamer).
%
% We detect Beamer via \@ifclassloaded{beamer}, which is the documented,
% non-internal way. This must be wrapped in \makeatletter/\makeatother
% because \@ifclassloaded contains an @-catcode character.
% -----------------------------------------------------------------------------
\makeatletter
\@ifclassloaded{beamer}{%
  \usefonttheme{professionalfonts}
  \setbeamercolor{structure}{fg=primary-blue}
  \setbeamercolor{frametitle}{fg=primary-blue}
  \setbeamercolor{title}{fg=primary-blue}
  \setbeamercolor{subtitle}{fg=primary-gold}
  \setbeamercolor{author}{fg=primary-blue}
  \setbeamercolor{institute}{fg=neutral}
  \setbeamercolor{date}{fg=primary-gold}
  \setbeamercolor{itemize item}{fg=primary-blue}
  \setbeamercolor{itemize subitem}{fg=primary-gold}
  \setbeamercolor{alerted text}{fg=primary-gold}
  \setbeamercolor{block title}{fg=white, bg=primary-blue}
  \setbeamercolor{block body}{bg=light-bg}
  % Semantic block variants: block=definition/concept (blue), exampleblock=
  % worked example (green/positive), alertblock=key takeaway or caution (gold).
  % Keep to <=2 colored boxes per slide across ALL three variants (INV-7).
  \setbeamercolor{block title example}{fg=white, bg=positive}
  \setbeamercolor{block body example}{bg=light-bg}
  \setbeamercolor{block title alerted}{fg=white, bg=primary-gold}
  \setbeamercolor{block body alerted}{bg=light-bg}

  % Minimal footer: course code (left) + page X / N (right), no navigation clutter
  \setbeamertemplate{navigation symbols}{}
  \setbeamertemplate{footline}{%
    \color{neutral}\footnotesize\hspace{0.7em}\@coursecodetext\hfill
    \insertframenumber/\inserttotalframenumber\hspace{0.7em}\vspace{0.3em}}
}{}
\makeatother

% -----------------------------------------------------------------------------
% TikZ setup — libraries the snippet gallery uses
% -----------------------------------------------------------------------------
\usepackage{tikz}
\usetikzlibrary{arrows.meta, positioning, calc, decorations.pathreplacing,
                fit, shapes.geometric, backgrounds}

% Shared TikZ styles — snippets and hand-written diagrams can pick these up
% via `\begin{tikzpicture}[every node/.style={...}]` or by naming specific
% styles (e.g., `\node[dag-node] (X) at (0,0) {X};`).
\tikzset{
  % Node styles for causal DAGs and similar
  dag-node/.style = {
    draw=primary-blue, fill=light-bg, rounded corners=2pt,
    minimum width=1.2cm, minimum height=0.9cm, align=center, font=\small
  },
  decision-node/.style = {
    draw=primary-gold, fill=white, diamond, aspect=1.8,
    minimum width=1.8cm, minimum height=1.2cm, align=center, font=\small,
    inner sep=1pt
  },
  % Edge styles — observed vs counterfactual
  observed-edge/.style = {-{Latex[length=2.5mm]}, thick, draw=jet},
  counterfactual-edge/.style = {-{Latex[length=2.5mm]}, thick, draw=neutral, dashed},
  confound-edge/.style = {-{Latex[length=2.5mm]}, thick, draw=negative, dashed},
  % Data-point styles for plots
  observed-dot/.style = {fill=primary-blue, draw=primary-blue, circle, inner sep=1.2pt},
  counterfactual-dot/.style = {fill=white, draw=neutral, circle, inner sep=1.2pt}
}

% -----------------------------------------------------------------------------
% Convenience macros
% -----------------------------------------------------------------------------
\newcommand{\muted}[1]{\textcolor{neutral}{#1}}
\newcommand{\key}[1]{\textcolor{primary-gold}{\textbf{#1}}}
\newcommand{\good}[1]{\textcolor{positive}{#1}}
\newcommand{\bad}[1]{\textcolor{negative}{#1}}

% Standout frame macro: full-bleed dark background for section transitions.
% Beamer-only (uses \begin{frame}/\setbeamercolor/\usebeamerfont) -- do not
% call from an article-class file (e.g. Notes/<CODE>/*-notes.tex). \muted,
% \key, \good, \bad, and \coursecode above are plain xcolor/gdef and are
% safe to use from either class; verified when Notes/article-class support
% was added.
% Expand: \transitionslide{Your transition title}
\newcommand{\transitionslide}[1]{%
  \begingroup
  \setbeamercolor{background canvas}{bg=primary-blue}
  \begin{frame}[plain]
    \centering
    \vfill
    {\usebeamerfont{title}\color{white}\Large #1\par}
    \vfill
  \end{frame}
  \endgroup
}

% Section divider frame (Beamer-only), an alternative to \transitionslide with a
% darker two-line style: near-black (jet) background, a small white label line
% above a larger gold title, and a thin gold rule along the bottom edge.
% Expand: \sectiondivider{Section 2}{Pointers}
\newcommand{\sectiondivider}[2]{%
  \begingroup
  \setbeamercolor{background canvas}{bg=jet}
  \begin{frame}[plain]
    \vspace*{3.0cm}
    {\color{white}\ttfamily\large #1\par}
    \vspace{0.5cm}
    {\color{primary-gold}\LARGE #2\par}
    \begin{tikzpicture}[remember picture, overlay]
      \draw[primary-gold, line width=2.5pt]
        ($(current page.south west)+(0,0.1)$) -- ($(current page.south east)+(0,0.1)$);
    \end{tikzpicture}
  \end{frame}
  \endgroup
}

% -----------------------------------------------------------------------------
% Channel watermark — "Concepts That Click" (YouTube).
%
% Article-class files (CompetitiveExam Q/A sets, Notes, Assignments) get a
% light diagonal draft watermark on every page plus a centered footer naming
% the channel. Beamer decks get a subtle TikZ background watermark instead
% (the footline already carries the course code). Centralized here so every
% MCQ PDF is branded without per-file edits.
% -----------------------------------------------------------------------------
\makeatletter
\@ifclassloaded{article}{%
  \usepackage{draftwatermark}
  \SetWatermarkText{Concepts That Click}
  \SetWatermarkScale{1}
  \SetWatermarkFontSize{2cm}
  \SetWatermarkLightness{0.86}
  \SetWatermarkAngle{45}
  \usepackage{fancyhdr}
  \pagestyle{fancy}
  \fancyhf{}
  \fancyfoot[C]{\footnotesize\textcolor{neutral}{Concepts That Click~|~YouTube: \href{https://www.youtube.com/@concepts.thatclick}{@concepts.thatclick}}}
  \renewcommand{\headrulewidth}{0pt}
  \renewcommand{\footrulewidth}{0pt}
  \fancypagestyle{plain}{%
    \fancyhf{}%
    \fancyfoot[C]{\footnotesize\textcolor{neutral}{Concepts That Click~|~YouTube: \href{https://www.youtube.com/@concepts.thatclick}{@concepts.thatclick}}}%
    \renewcommand{\headrulewidth}{0pt}%
    \renewcommand{\footrulewidth}{0pt}%
  }
}{}
\@ifclassloaded{beamer}{%
  \setbeamertemplate{background}{%
    \begin{tikzpicture}[remember picture, overlay]
      \node[rotate=45, opacity=0.055, align=center]
        at (current page.center)
        {\Huge\textcolor{neutral}{Concepts That Click}};
    \end{tikzpicture}%
  }
}{}
\makeatother 
\coursecode{JKSSB}
\title{Lesson 50 Practice: Data Representation, Number Systems \& Logic}
\author{JKSSB Computer Awareness}
\date{\today}
\begin{document}
\maketitle
\noindent\textbf{Provenance:} 16 original JKSSB-pattern questions and five
paraphrased adaptations of official JKSSB PYQs: Computer Instructor/Operator
2025 Q1--Q4 and Website Operator 2026 Q73. Adaptations are not verbatim
reproductions.

\begin{enumerate}[label=\textbf{Q\arabic*.}, leftmargin=*]
\item \textit{[PYQ-adapted: JKSSB Computer Instructor/Operator 2025 Q1;
official paper; paraphrased]} What is the decimal ASCII code for the dollar
sign (\$)?
  \begin{enumerate}[label=\Alph*.]
    \item 32
    \item 30
    \item 34
    \item 36
  \end{enumerate}
\item \textit{[PYQ-adapted: JKSSB Computer Instructor/Operator 2025 Q2;
official paper; paraphrased]} Convert the binary number
$11010101.01_2$ to decimal.
  \begin{enumerate}[label=\Alph*.]
    \item $200.25_{10}$
    \item $198.25_{10}$
    \item $213.25_{10}$
    \item $209.25_{10}$
  \end{enumerate}
\item \textit{[PYQ-adapted: JKSSB Computer Instructor/Operator 2025 Q3;
official paper; paraphrased]} Which statement about UTF-8 is correct?
  \begin{enumerate}[label=\Alph*.]
    \item It is fixed-length and uses exactly 8 bits for every character.
    \item It is variable-length and compatible with ASCII in the ASCII range.
    \item It assigns Unicode code points only to Latin letters.
    \item It stores every character using 8 bytes.
  \end{enumerate}
\item \textit{[PYQ-adapted: JKSSB Computer Instructor/Operator 2025 Q4;
official paper; paraphrased]} What is the binary representation of
$278_{10}$?
  \begin{enumerate}[label=\Alph*.]
    \item $100010110_2$
    \item $101010110_2$
    \item $100110110_2$
    \item $110010110_2$
  \end{enumerate}
\item \textit{[PYQ-adapted: JKSSB Website Operator 2026 Q73; official
provisional key; paraphrased]} Which binary numeral has the value
$10_{10}$?
  \begin{enumerate}[label=\Alph*.]
    \item $1010_2$
    \item $1001_2$
    \item $1100_2$
    \item $1110_2$
  \end{enumerate}
\item \textit{[Original]} What is the base of the octal number system?
  \begin{enumerate}[label=\Alph*.]
    \item 2
    \item 8
    \item 10
    \item 16
  \end{enumerate}
\item \textit{[Original]} Which digit is not permitted in an octal numeral?
  \begin{enumerate}[label=\Alph*.]
    \item 0
    \item 5
    \item 7
    \item 8
  \end{enumerate}
\item \textit{[Original]} What is the decimal value of $2F_{16}$?
  \begin{enumerate}[label=\Alph*.]
    \item 31
    \item 32
    \item 47
    \item 63
  \end{enumerate}
\item \textit{[Original]} Convert $1111111_2$ to octal.
  \begin{enumerate}[label=\Alph*.]
    \item $177_8$
    \item $177_2$
    \item $77_8$
    \item $177_{16}$
  \end{enumerate}
\item \textit{[Original]} In BCD, how is the decimal number 25 represented?
  \begin{enumerate}[label=\Alph*.]
    \item $11001_2$
    \item $0010\ 0101$
    \item $0101\ 0010$
    \item $11101_2$
  \end{enumerate}
\item \textit{[Original]} Which statement is \textbf{NOT} correct?
  \begin{enumerate}[label=\Alph*.]
    \item Original ASCII uses 7-bit codes.
    \item Unicode assigns code points across many writing systems.
    \item UTF-8 is variable-length.
    \item BCD and Unicode are interchangeable names for character encoding.
  \end{enumerate}
\item \textit{[Original]} Which option describes the relationship between
  Unicode and UTF-8 correctly?
  \begin{enumerate}[label=\Alph*.]
    \item Unicode assigns code points; UTF-8 encodes them using a
      variable-length byte sequence.
    \item UTF-8 is another name for BCD.
    \item Unicode is a 4-bit code for each decimal digit.
    \item UTF-8 is always fixed at 8 bytes per character.
  \end{enumerate}
\item \textit{[Original]} In common 24-bit RGB, the three 8-bit channels
  represent:
  \begin{enumerate}[label=\Alph*.]
    \item Red, green, and blue
    \item Red, grey, and black
    \item Row, grid, and bit
    \item Read, generate, and store
  \end{enumerate}
\item \textit{[Original]} For inputs $A=1$ and $B=0$, what is the output of
  XOR?
  \begin{enumerate}[label=\Alph*.]
    \item 0
    \item 1
    \item 2
    \item It depends on a clock
  \end{enumerate}
\item \textit{[Original]} Which gate outputs 1 only when both inputs are 1?
  \begin{enumerate}[label=\Alph*.]
    \item OR
    \item XOR
    \item AND
    \item NOT
  \end{enumerate}
\item \textit{[Original]} For $A=1$ and $B=1$, the NAND output is:
  \begin{enumerate}[label=\Alph*.]
    \item 0
    \item 1
    \item 2
    \item Undefined
  \end{enumerate}
\item \textit{[Original]} A NOT gate has:
  \begin{enumerate}[label=\Alph*.]
    \item One input and the opposite output
    \item Two inputs and output 1 only when both are 1
    \item Two inputs and output 1 when they differ
    \item Three inputs representing RGB
  \end{enumerate}
\item \textit{[Original]} Evaluate the statements about a common 24-bit RGB
  representation:
  \begin{enumerate}[label=\Roman*.]
    \item It uses red, green, and blue channels.
    \item Three 8-bit channels total 24 bits per pixel.
    \item Every image format must use exactly 24-bit RGB.
  \end{enumerate}
  \begin{enumerate}[label=\Alph*.]
    \item I only
    \item I and II only
    \item II and III only
    \item I, II, and III
  \end{enumerate}
\item \textit{[Original]} What is $1011_2$ in decimal?
  \begin{enumerate}[label=\Alph*.]
    \item 9
    \item 10
    \item 11
    \item 13
  \end{enumerate}
\item \textit{[Original]} When converting a binary numeral to hexadecimal by
  grouping bits, how many bits are placed in each complete group?
  \begin{enumerate}[label=\Alph*.]
    \item 2
    \item 3
    \item 4
    \item 8
  \end{enumerate}
\item \textit{[Original]} For inputs $A=0$ and $B=0$, what are the outputs
  of AND, OR, XOR, and NAND, in that order?
  \begin{enumerate}[label=\Alph*.]
    \item $0,0,0,1$
    \item $1,0,0,0$
    \item $0,1,1,1$
    \item $1,1,0,1$
  \end{enumerate}
\end{enumerate}
\end{document}
